% =====================================================================
%  Protocolo de survey pré-registrado — Parte 1 (Planejamento)
%  Medição e Experimentação em Engenharia de Software — 2º/2026
%  Compilar com:  tectonic protocolo.tex   (ou xelatex/lualatex)
%  Marcadores \pend{...} = campos que o grupo DEVE preencher antes do envio.
% =====================================================================
\documentclass[10pt,a4paper]{article}

\usepackage[a4paper,top=1.6cm,bottom=1.7cm,left=1.8cm,right=1.8cm,
            headheight=13pt,headsep=7pt,footskip=18pt]{geometry}
\usepackage[brazilian]{babel}
\usepackage{libertinus}
\usepackage{microtype}
\usepackage[table]{xcolor}
\usepackage{booktabs,tabularx,longtable,array,multirow,float,needspace}
\usepackage{enumitem}
\usepackage{multicol}
\usepackage{titlesec}
\usepackage{fancyhdr}
\usepackage{lastpage}
\usepackage{amsmath}
\usepackage{caption}
\usepackage[authoryear,round]{natbib}
\setcitestyle{aysep={,},citesep={;}}
\usepackage[most]{tcolorbox}
\usepackage[hidelinks]{hyperref}

% ---------- Paleta ----------------------------------------------------
% Somente preto e branco. O laranja de \pend{} marca apenas campos provisórios.
\definecolor{pendc}{HTML}{B54708}

% Campo pendente (deve ser preenchido pelo grupo antes do envio)
\newcommand{\pend}[1]{\textcolor{pendc}{\textbf{[}\,#1\,\textbf{]}}}

% ---------- Tipografia de títulos -------------------------------------
\titleformat{\section}{\normalfont\bfseries\large}
  {\thesection}{0.6em}{}
\titlespacing*{\section}{0pt}{1.3ex plus .3ex}{0.6ex}
\titleformat{\subsection}[runin]{\normalfont\bfseries\normalsize}
  {\thesubsection}{0.45em}{}[.]
\titlespacing*{\subsection}{0pt}{0.9ex plus .2ex}{0.5em}

% ---------- Cabeçalho e rodapé ----------------------------------------
\pagestyle{fancy}\fancyhf{}
\renewcommand{\headrulewidth}{0.4pt}
\fancyhead[L]{\footnotesize\itshape Protocolo de survey pré-registrado · Dificuldade no aprendizado de programação}
\fancyhead[R]{\footnotesize\itshape MEES · 2º/2026}
\fancyfoot[L]{\footnotesize Versão 1.0 · registrada em 13/11/2026}
\fancyfoot[R]{\footnotesize\thepage\ de \pageref{LastPage}}
\fancypagestyle{first}{\fancyhf{}\renewcommand{\headrulewidth}{0pt}
  \fancyfoot[L]{\footnotesize Versão 1.0 · registrada em 13/11/2026}
  \fancyfoot[R]{\footnotesize\thepage\ de \pageref{LastPage}}}

% ---------- Tabelas, listas, legendas ---------------------------------
\captionsetup{font=small,labelfont=bf,skip=3pt,justification=raggedright,singlelinecheck=false}
\captionsetup[table]{position=top}
\renewcommand{\arraystretch}{1.06}
\setlength{\tabcolsep}{4pt}
\setlist{nosep,leftmargin=1.2em}
\newcolumntype{L}[1]{>{\raggedright\arraybackslash}p{#1}}
\newcolumntype{Y}{>{\raggedright\arraybackslash}X}
\newcommand{\thead}[1]{\textbf{#1}}
\newcommand{\blockrow}[2]{\multicolumn{#1}{@{}l}{\textbf{\textit{#2}}}\\}
\setlength{\parskip}{0.25ex}
\setlength{\parindent}{0pt}
\setlength{\LTpre}{9pt}\setlength{\LTpost}{2pt}
\hypersetup{pdftitle={Protocolo de survey pré-registrado: dificuldade percebida no aprendizado de programação},
            pdfauthor={Grupo — Medição e Experimentação em ES}}

% =====================================================================
\begin{document}
\thispagestyle{first}


% ---------- Bloco de título -------------------------------------------
\begin{center}
{\footnotesize\textsc{Protocolo de survey pré-registrado} \,·\, \textsc{Parte 1 — Planejamento}}\\[0.5ex]
{\LARGE\bfseries Fatores associados à dificuldade percebida no\\[0.3ex] aprendizado de programação}\\[0.5ex]
{\normalsize\itshape Survey transversal com estudantes de cursos de tecnologia: períodos iniciais e estagiários}\\[0.8ex]
{\small Ana Luiza Santos Gomes, Bruna Barbosa Portilho Bernardes, Guilherme Augustto Costa Barros e Vinícius Dias}\\[0.4ex]
{\small \pend{Nome da instituição} \,·\, Medição e Experimentação em Engenharia de Software (MEES) \,·\, 2º/2026}\\[0.4ex]
{\small Registro: tag \texttt{protocolo-v1.0} em \href{https://github.com/guicostak/Trabalho-Final-Survey}{\texttt{github.com/guicostak/Trabalho-Final-Survey}}}
\end{center}

\vspace{0.8ex}
\begin{tcolorbox}[enhanced,colback=white,colframe=black,boxrule=0.5pt,arc=0pt,
  left=6pt,right=6pt,top=4pt,bottom=4pt,fontupper=\small]
\textbf{Ficha do estudo.}
\textbf{Método:} survey por questionário autoadministrado on-line, descritivo e correlacional, de corte transversal \textbullet\
\textbf{População:} estudantes de CC, ES, ADS e SI de instituições brasileiras, do 1º ao 4º período ou em estágio de TI \textbullet\
\textbf{Amostragem:} não probabilística, por conveniência com bola de neve, em vários canais e com metas por segmento \textbullet\
\textbf{Amostra:} alvo de 100 respostas válidas (mínimo de 50); margem de erro nominal de ±9,8 p.p. \textbullet\
\textbf{Hipóteses confirmatórias:} 8, testadas com $\alpha=0{,}05$ e correção de Holm \textbullet\
\textbf{Coleta:} 18/11 a 04/12/2026 \textbullet\
\textbf{Materiais:} instrumento completo, livro de códigos e TCLE no repositório do grupo.
O protocolo segue o padrão \emph{Questionnaire Surveys} e o suplemento \emph{Registered Reports} dos ACM SIGSOFT Empirical Standards \citep{ralph} e foi registrado \emph{antes} da primeira resposta. Qualquer desvio será datado, justificado e registrado em \texttt{DESVIOS.md}.
\end{tcolorbox}

% =====================================================================
\Needspace{5\baselineskip}\section{Objetivo e questões de pesquisa}

Disciplinas introdutórias de programação têm reprovação elevada: a taxa média de aprovação fica próxima de dois terços \citep{watson}. A literatura relaciona a dificuldade de aprender a programar a fatores de \emph{entrada} (experiência prévia), \emph{psicológicos} (autoeficácia e ansiedade) e de \emph{conteúdo} (tópicos como recursão e ponteiros) \citep{robins,lahtinen,ramalingam,connolly}. Além disso, os assistentes de IA generativa mudaram a forma como os iniciantes estudam \citep{becker}. Este estudo caracteriza a dificuldade \emph{percebida} (a percepção do estudante, não o seu desempenho) em quatro cursos de tecnologia e compara estudantes dos períodos iniciais com estagiários, que já aplicam programação em contexto profissional. O objetivo segue o template GQM \citep{basili} (Tabela~\ref{tab:gqm}); as questões de pesquisa (QP) e as hipóteses estão na Tabela~\ref{tab:qp}, e as métricas na Tabela~\ref{tab:met}.

\begin{table}[H]
\centering\small
\caption{Objetivo de medição no template GQM.}\label{tab:gqm}
\begin{tabularx}{\textwidth}{@{}L{3.1cm}Y@{}}
\toprule
\thead{Analisar} & o processo de aprendizagem de programação na graduação \\
\thead{com o propósito de} & caracterizar e identificar fatores associados \\
\thead{com respeito à} & dificuldade percebida pelos estudantes \\
\thead{do ponto de vista} & dos estudantes de graduação \\
\thead{no contexto de} & cursos de Ciência da Computação (CC), Engenharia de Software (ES), Análise e Desenvolvimento de Sistemas (ADS) e Sistemas de Informação (SI) de instituições brasileiras, com estudantes do 1º ao 4º período ou em estágio de TI, no 2º semestre de 2026 \\
\bottomrule
\end{tabularx}
\end{table}

\vspace{-1ex}
\begin{table}[H]
\centering\small
\caption{Questões de pesquisa e hipóteses (bilaterais; direção esperada entre parênteses).}\label{tab:qp}
\begin{tabularx}{\textwidth}{@{}L{0.75cm}L{5.9cm}Y@{}}
\toprule
\thead{QP} & \thead{Pergunta} & \thead{Hipóteses} \\
\midrule
QP1 & Qual é o nível de dificuldade percebida no aprendizado de programação e quais tópicos são percebidos como mais difíceis? & Descritiva, sem hipótese. \\
QP2 & A dificuldade percebida difere entre estudantes com e sem experiência prévia em programação? & H2a: M1 difere entre os grupos (maior sem experiência).\newline H2b: reprovação prévia está associada à experiência prévia. \\
QP3 & A dificuldade percebida está associada à autoeficácia e à ansiedade em programação? & H3a: $\rho(\text{M1},\text{M4})\neq0$ (negativa).\newline H3b: $\rho(\text{M1},\text{M5})\neq0$ (positiva). \\
QP4 & A dificuldade percebida varia conforme o curso, a realização de estágio e os hábitos de estudo? & H4a: M1 difere entre os cursos.\newline H4b: M1 difere entre estagiários e não estagiários (menor entre estagiários).\newline H4c: $\rho(\text{M1},\text{M6})\neq0$.\newline H4d: $\rho(\text{M1},\text{M7})\neq0$ (sem direção esperada). \\
\bottomrule
\end{tabularx}
\end{table}

{\small
\begin{longtable}{@{}L{0.65cm}L{3.0cm}L{8.55cm}L{2.45cm}L{1.55cm}@{}}
\caption{Métricas: definição operacional, escala de medição e QP atendida.}\label{tab:met}\\
\toprule
\thead{ID} & \thead{Métrica} & \thead{Definição operacional} & \thead{Escala} & \thead{QP} \\
\midrule\endfirsthead
\toprule
\thead{ID} & \thead{Métrica} & \thead{Definição operacional} & \thead{Escala} & \thead{QP} \\
\midrule\endhead
M1 & Dificuldade percebida & Média de DIF1–DIF4 (DIF4 invertido), de 1 a 5 & Ordinal (escore)\textsuperscript{a} & QP1–QP4 \\
M2 & Dificuldade por tópico & \% de respostas \emph{difícil} ou \emph{muito difícil} em cada TOP1–TOP8 & Razão (proporção) & QP1 \\
M3 & Experiência prévia & Programava antes do ingresso (PER7); tempo (PER8); ensino técnico (PER9) & Nominal; ordinal & QP2 \\
M4 & Autoeficácia & Média de AUT1–AUT4, de 1 a 5 & Ordinal (escore)\textsuperscript{a} & QP3 \\
M5 & Ansiedade em programação & Média de ANS1–ANS3, de 1 a 5 & Ordinal (escore)\textsuperscript{a} & QP3 \\
M6 & Prática semanal & Horas semanais de programação fora da aula (HAB1) & Razão & QP4 \\
M7 & Uso de IA generativa & Frequência de uso de IA ao estudar (HAB2), de 1 a 5 & Ordinal & QP4 \\
M8 & Período & Período em curso (PER2) & Ordinal & QP4 (expl.) \\
M9 & Reprovação prévia & Já teve reprovação em disciplina de programação (PER10) & Nominal & QP2 \\
M10 & Curso & CC, ES, ADS ou SI (PER1) & Nominal & QP4 \\
M11 & Estágio & Faz estágio em TI atualmente (PER3) & Nominal & QP4 \\
M12 & Composição da amostra & Tipo de instituição (PER4), canal de convite (PER5) e faixa etária (PER6) & Nominal; ordinal & QP1–QP4\textsuperscript{b} \\
M13 & Qualidade das respostas & \% de respostas válidas e descartes por motivo (TRI1–TRI3, ATN1; regras E1–E8) & Razão (proporção) & QP1–QP4\textsuperscript{b} \\
M14 & Modo de uso da IA & Principal uso da IA ao estudar (HAB3) & Nominal & QP4 \\
M15 & Temas da questão aberta & Frequência de cada tema codificado em ABR1 & Nominal & QP1 \\
\bottomrule
\multicolumn{5}{@{}l}{\footnotesize\textsuperscript{a}Itens Likert são ordinais; o escore médio entra no $\alpha$ de Cronbach, mas só é analisado com testes de postos. \textsuperscript{b}Validade e generalização de todas as QPs.}
\end{longtable}
}

% =====================================================================
\Needspace{5\baselineskip}\section{População e amostragem}

\subsection{População-alvo}
São estudantes de graduação em CC, ES, ADS e SI de instituições brasileiras, públicas ou privadas, com 18 anos ou mais, que estejam (a) do 1º ao 4º período ou (b) em qualquer período, fazendo estágio em TI. Estudantes formam a população de interesse, e não um substituto por conveniência: a dificuldade de aprender a programar se manifesta nos períodos iniciais, e os estagiários permitem contrastar quem já aplica programação em contexto profissional. Não existe lista nacional desses estudantes, e o tamanho da população é estimado em $N\approx$ \pend{N}. Ficam de fora os menores de 18 anos, que exigiriam autorização do responsável. Os colegas da disciplina MEES e os integrantes do grupo não recebem o convite e são orientados a não responder, porque conhecem as hipóteses do estudo.

\subsection{Origem dos contatos e estratégia}
Como não há quadro amostral, a amostragem é \emph{não probabilística, por conveniência com bola de neve} \citep{baltes}. O convite parte de cinco fontes:
(i) docentes de programação de outras instituições, contatados pelo grupo, que repassam o convite às turmas;
(ii) centros e diretórios acadêmicos;
(iii) grupos de estudantes de tecnologia (WhatsApp, Discord, Telegram), com permissão dos administradores;
(iv) perfis dos integrantes no LinkedIn e no Instagram;
(v) os próprios respondentes, convidados a repassar o link.
O formulário registra o canal de chegada (PER5) e o tipo de instituição (PER4), mas não o nome da instituição, o que reduz o risco de reidentificação. Para que as comparações sejam viáveis, há \textbf{metas mínimas por segmento}: 15 respostas por curso e 20 estagiários. Em 25/11, a composição parcial é verificada, e a divulgação é direcionada aos segmentos abaixo da meta. A elegibilidade é verificada no próprio formulário (Bloco 0 e itens PER1–PER3).

\subsection{Tamanho da amostra, margem de erro e poder}
O alvo é de \textbf{100 respostas válidas}, acima do mínimo de 50 exigido. A margem de erro de uma proporção com 95\% de confiança é $e = 1{,}96\sqrt{p(1-p)/n}$, no pior caso $p=0{,}5$, sem correção para população finita, porque a população é grande e de tamanho desconhecido. Como a amostra não é probabilística, essa margem é \textbf{nominal}: indica a precisão que uma amostra aleatória do mesmo tamanho teria e não garante representatividade (Seção~7). O poder estatístico é de 80\% com $\alpha=0{,}05$ bilateral (Tabela~\ref{tab:me}).

\begin{table}[H]
\centering\small
\caption{Margem de erro nominal (95\%) e menores efeitos detectáveis com poder de 80\%, por tamanho de amostra.}\label{tab:me}
\begin{tabularx}{\textwidth}{@{}L{3.2cm}*{3}{>{\centering\arraybackslash}X}@{}}
\toprule
\thead{Respostas válidas ($n$)} & \thead{Margem de erro nominal} & \thead{Menor $|\rho|$ detectável} & \thead{Menor $d$ detectável\textsuperscript{c}} \\
\midrule
50 (mínimo) & ±13,9 p.p. & 0,39 & 0,81 \\
\textbf{100 (alvo)} & \textbf{±9,8 p.p.} & \textbf{0,28} & \textbf{0,57} \\
150 & ±8,0 p.p. & 0,23 & 0,47 \\
\bottomrule
\multicolumn{4}{@{}l}{\footnotesize\textsuperscript{c}Mann-Whitney com grupos de mesmo tamanho (eficiência relativa assintótica de 0,955 em relação ao teste $t$).}
\end{tabularx}
\end{table}

% =====================================================================
\Needspace{5\baselineskip}\section{Questionário com rastreabilidade}

O instrumento tem 37 itens em oito blocos e deve levar cerca de 9 minutos. Ele será aplicado no Microsoft Forms, aberto a qualquer pessoa com o link e configurado para \emph{não registrar nome nem e-mail}. A plataforma registra o início e o fim de cada resposta, o que permite detectar respostas rápidas demais. A ordem dos itens é fixa. Todas as escalas Likert têm 5 pontos rotulados (1 = \emph{discordo totalmente} … 5 = \emph{concordo totalmente}), com ponto neutro. O instrumento segue as recomendações de Kitchenham e Pfleeger (2008) para a redação de itens. Os itens de autoeficácia foram adaptados de Ramalingam e Wiedenbeck (1998) e de Tsai, Wang e Hsu (2018); os de ansiedade, de Connolly, Murphy e Moore (2009); e a lista de tópicos, de Lahtinen, Ala-Mutka e Järvinen (2005). Os itens de dificuldade percebida foram elaborados pelo grupo, evitando sobreposição com os de autoeficácia, e terão o conteúdo validado no pré-teste (Seção~4). As instruções pedem que o estudante considere todo o seu aprendizado de programação até agora. Cada construto é medido por 3 ou 4 itens, e DIF4 é invertido para reduzir a aquiescência. O item ATN1 é a \textbf{pergunta de atenção}. A redação completa, as opções e as regras de ramificação estão em \texttt{instrumento/questionario.md}. A Tabela~\ref{tab:inst} liga cada item a uma escala, a uma métrica e a uma QP.

{\footnotesize
\begin{longtable}{@{}L{1.05cm}L{8.75cm}L{2.45cm}L{1.15cm}L{1.95cm}@{}}
\caption{Matriz de rastreabilidade do questionário (item → escala → métrica → QP). \textsuperscript{R} = item de codificação invertida.}\label{tab:inst}\\
\toprule
\thead{ID} & \thead{Item (redação resumida)} & \thead{Escala} & \thead{Métrica} & \thead{QP} \\
\midrule\endfirsthead
\toprule
\thead{ID} & \thead{Item (redação resumida)} & \thead{Escala} & \thead{Métrica} & \thead{QP} \\
\midrule\endhead
\blockrow{5}{Bloco 0 — Consentimento e triagem (respostas que indicam inelegibilidade encerram o formulário)}
TRI1 & Li o Termo de Consentimento e concordo em participar & Nominal & M13 & QP1–QP4 \\
TRI2 & Tenho 18 anos ou mais & Nominal & M13 & QP1–QP4 \\
TRI3 & Já respondi a este questionário antes & Nominal & M13 & QP1–QP4 \\
\blockrow{5}{Bloco A — Perfil e elegibilidade}
PER1 & Curso (CC, ES, ADS, SI ou outro; \emph{outro} encerra) & Nominal & M10 & QP4 \\
PER2 & Período em curso (1º, 2º, 3º, 4º, 5º ou posterior) & Ordinal & M8 & QP4 (expl.) \\
PER3 & Faz estágio em TI atualmente (5º período ou posterior sem estágio encerra) & Nominal & M11 & QP4 \\
PER4 & Tipo de instituição (pública ou privada) & Nominal & M12 & QP1–QP4 \\
PER5 & Canal pelo qual recebeu o convite & Nominal & M12 & QP1–QP4 \\
PER6 & Faixa etária (18–19, 20–22, 23–25, 26–30, 31+) & Ordinal & M12 & QP1–QP4 \\
PER7 & Programava antes de ingressar no curso & Nominal & M3 & QP2 \\
PER8 & Tempo de experiência prévia (<6 meses … >2 anos; condicional a PER7) & Ordinal & M3 & QP2 \\
PER9 & Cursou ensino médio técnico em informática ou área afim & Nominal & M3 & QP2 \\
PER10 & Já teve reprovação em disciplina de programação no curso (opção \emph{prefiro não responder}) & Nominal & M9 & QP2 \\
\blockrow{5}{Bloco B — Dificuldade percebida (construto DIF)}
DIF1 & Tenho dificuldade para compreender novos conceitos de programação & Ordinal (Likert 5) & M1 & QP1–QP4 \\
DIF2 & Tenho dificuldade para transformar a lógica de uma solução em código & Ordinal (Likert 5) & M1 & QP1–QP4 \\
DIF3 & Com frequência não sei por onde começar a resolver um problema de programação & Ordinal (Likert 5) & M1 & QP1–QP4 \\
DIF4\textsuperscript{R} & Consigo acompanhar o ritmo das disciplinas de programação sem grandes dificuldades & Ordinal (Likert 5) & M1 & QP1–QP4 \\
\blockrow{5}{Bloco C — Dificuldade por tópico (1 = muito fácil … 5 = muito difícil; \emph{ainda não estudei} = ausente)}
TOP1–2 & Estruturas condicionais; estruturas de repetição & Ordinal (5 pts) & M2 & QP1 \\
TOP3–4 & Vetores e matrizes; funções/métodos e passagem de parâmetros & Ordinal (5 pts) & M2 & QP1 \\
TOP5–6 & Recursão; ponteiros, referências e gerenciamento de memória & Ordinal (5 pts) & M2 & QP1 \\
TOP7–8 & Depuração (localizar e corrigir erros); orientação a objetos & Ordinal (5 pts) & M2 & QP1 \\
\blockrow{5}{Bloco D — Autoeficácia em programação (construto AUT) e verificação de atenção}
AUT1 & Consigo escrever um programa que resolva um problema simples proposto em aula & Ordinal (Likert 5) & M4 & QP3 \\
AUT2 & Consigo localizar e corrigir erros em um programa de minha autoria & Ordinal (Likert 5) & M4 & QP3 \\
AUT3 & Consigo entender um programa escrito por outra pessoa, de complexidade compatível com a disciplina & Ordinal (Likert 5) & M4 & QP3 \\
AUT4 & Consigo concluir um exercício sem ajuda de colegas, monitores ou ferramentas de IA & Ordinal (Likert 5) & M4 & QP3 \\
ATN1 & \textbf{Atenção:} ``Para confirmar que está lendo, marque \emph{Discordo totalmente} nesta linha'' & Nominal (acerto) & M13 & QP1–QP4 \\
\blockrow{5}{Bloco E — Ansiedade em programação (construto ANS)}
ANS1 & Sinto nervosismo quando preciso programar em provas ou atividades avaliativas & Ordinal (Likert 5) & M5 & QP3 \\
ANS2 & Sinto desânimo quando meu código apresenta erros que não consigo resolver & Ordinal (Likert 5) & M5 & QP3 \\
ANS3 & Evito iniciar tarefas de programação por receio de não conseguir concluí-las & Ordinal (Likert 5) & M5 & QP3 \\
\blockrow{5}{Bloco F — Hábitos de estudo e bloco G — questão aberta}
HAB1 & Horas por semana dedicadas a programar fora da aula (inteiro, 0–40) & Razão & M6 & QP4 \\
HAB2 & Frequência de uso de IA generativa ao estudar (nunca … sempre) & Ordinal (5 pts) & M7 & QP4 \\
HAB3 & Principal uso da IA (código pronto, entender conceitos/erros, revisar código, outro, não uso) & Nominal & M14 & QP4 \\
ABR1 & O que mais dificulta ou mais ajudou o seu aprendizado de programação? (opcional) & Texto → nominal & M15 & QP1 \\
\bottomrule
\end{longtable}
}

% =====================================================================
\Needspace{5\baselineskip}\section{Pré-teste}

O pré-teste será aplicado em 13/11/2026 a \pend{4} pessoas que \emph{não} fazem parte da amostra: \pend{3} estudantes do 5º período ou posterior que não fazem estágio, e por isso estão fora da população, de preferência de cursos diferentes; e \pend{1} docente de programação introdutória, que fará a revisão de conteúdo por especialista. Cada participante responderá ao formulário final pensando em voz alta (entrevista cognitiva). O grupo cronometrará a duração e fará uma conversa de encerramento com quatro perguntas: clareza dos itens, ambiguidade, duração percebida e itens sensíveis ou faltantes. Um item é reformulado quando ao menos dois participantes ou o especialista apontam um problema nele. As respostas do pré-teste são descartadas e não entram no conjunto de dados. O roteiro e o registro detalhado estão em \texttt{pre-teste/roteiro-pre-teste.md}.

\begin{table}[H]
\centering\small
\caption{Resultado do pré-teste e mudanças feitas no instrumento.}\label{tab:pre}
\begin{tabularx}{\textwidth}{@{}L{1.4cm}L{4.2cm}YL{1.6cm}@{}}
\toprule
\thead{Item} & \thead{Problema observado (quem)} & \thead{Mudança realizada} & \thead{Duração} \\
\midrule
\pend{ID} & \pend{descrição; P1, P3} & \pend{nova redação ou remoção} & \pend{min} \\
\pend{ID} & \pend{descrição; especialista} & \pend{nova redação ou remoção} & \pend{min} \\
\pend{ID} & \pend{descrição} & \pend{nova redação ou remoção} & \pend{min} \\
\bottomrule
\end{tabularx}
\end{table}
\vspace{-1ex}
A duração mediana observada no pré-teste foi de \pend{X} min. O valor de referência usado na regra E7 (Seção~5) é a mediana da amostra real.

% =====================================================================
\Needspace{5\baselineskip}\section{Plano de análise}

\subsection{Limpeza}
As regras da Tabela~\ref{tab:limpeza} são aplicadas nesta ordem, e cada descarte é registrado com o seu motivo para alimentar o perfil da amostra no dashboard. Como o link é aberto, a taxa de resposta não pode ser calculada, o que os Empirical Standards aceitam para surveys abertos. Em seu lugar, relatam-se as respostas por canal (PER5) e por tipo de instituição (PER4). Os registros descartados permanecem no arquivo bruto, para auditoria.

\begin{table}[H]
\centering\small
\caption{Regras de exclusão de respostas, na ordem de aplicação.}\label{tab:limpeza}
\begin{tabularx}{\textwidth}{@{}L{0.9cm}L{4.0cm}YL{2.2cm}@{}}
\toprule
\thead{Regra} & \thead{Motivo do descarte} & \thead{Critério operacional} & \thead{Itens} \\
\midrule
E1 & Sem consentimento & Não concordou com o TCLE & TRI1 \\
E2 & Menor de 18 anos & Declarou ter menos de 18 anos & TRI2 \\
E3 & Resposta repetida declarada & Declarou já ter respondido antes & TRI3 \\
E4 & Fora da população & Curso \emph{outro}, ou 5º período ou posterior sem estágio em TI & PER1–PER3 \\
E5 & Falha na atenção & Qualquer resposta diferente de \emph{Discordo totalmente} & ATN1 \\
E6 & Resposta retilínea & Variância zero nos 11 itens Likert, que têm sentidos opostos & DIF, AUT, ANS \\
E7 & Resposta rápida demais & Duração menor que 1/3 da mediana das durações & Início/fim \\
E8 & Duplicata técnica & Itens fechados idênticos enviados com menos de 5 min de diferença (mantém-se a primeira) & Todos \\
\bottomrule
\end{tabularx}
\end{table}

\vspace{-1ex}
\subsection{Codificação e dados ausentes}
Os itens Likert são codificados de 1 a 5, e DIF4 é invertido ($6-x$). Os escores M1, M4 e M5 são a média dos itens do construto. Como os itens fechados são obrigatórios, só há ausência por desenho: \emph{ainda não estudei} (TOP), \emph{prefiro não responder} (PER10) e ABR1, que é opcional. Esses casos são tratados por exclusão par a par, sem imputação. Valores de HAB1 acima de 40 são considerados implausíveis e tratados como ausentes.

\subsection{Confiabilidade}
Calcula-se o $\alpha$ de Cronbach de DIF, AUT e ANS, com a correlação item-total. Os critérios são: ${\alpha \geq 0{,}70}$ adequado; de 0,60 a 0,70 aceitável com ressalva; abaixo de 0,60 os itens passam a ser analisados separadamente, o que é declarado como desvio.

\subsection{Estatística descritiva e testes}
O relatório inclui frequências, mediana e intervalo interquartil, gráficos de barras divergentes para os itens Likert e proporções com IC de 95\% (Wilson) acompanhadas da margem de erro. Nenhum teste supõe normalidade. Os testes e os tamanhos de efeito estão na Tabela~\ref{tab:testes}.

\begin{table}[H]
\centering\small
\caption{Testes coerentes com as escalas e tamanhos de efeito por hipótese.}\label{tab:testes}
\begin{tabularx}{\textwidth}{@{}L{1.35cm}L{4.2cm}L{4.1cm}Y@{}}
\toprule
\thead{Hipótese} & \thead{Variáveis (escalas)} & \thead{Teste} & \thead{Tamanho de efeito (pequeno/médio/grande)} \\
\midrule
QP1 & M1, M2 (ordinal; proporção) & Descritiva; IC de 95\% de Wilson & Proporção ± margem de erro \\
H2a & M1 × M3 (ordinal × nominal, 2 grupos) & Mann-Whitney $U$ & $\delta$ de Cliff: 0,147/0,33/0,474 \citep{romano} \\
H2b & M3 × M9 (nominal × nominal) & $\chi^2$ de independência\textsuperscript{d} & V de Cramér: 0,1/0,3/0,5 \citep{cohen} \\
H3a, H3b & M1 × M4; M1 × M5 (ordinal × ordinal) & Correlação de Spearman & $\rho$ com IC bootstrap (2.000 reamostras): 0,1/0,3/0,5 \\
H4a & M1 × M10 (ordinal × 4 grupos) & Kruskal-Wallis; post hoc de Dunn & $\varepsilon^2$: 0,01/0,08/0,26 \citep{tomczak} \\
H4b & M1 × M11 (ordinal × nominal, 2 grupos) & Mann-Whitney $U$ & $\delta$ de Cliff: 0,147/0,33/0,474 \\
H4c, H4d & M1 × M6 (razão); M1 × M7 (ordinal) & Correlação de Spearman & $\rho$: 0,1/0,3/0,5 \\
\bottomrule
\multicolumn{4}{@{}l}{\footnotesize\textsuperscript{d}Teste exato de Fisher quando alguma frequência esperada for menor que 5.}
\end{tabularx}
\end{table}

\vspace{-1.5ex}
\subsection{Decisão e múltiplas comparações}
O nível de significância é $\alpha=0{,}05$, com testes bilaterais. A \textbf{correção de Holm} \citep{holm} é aplicada à família das oito hipóteses confirmatórias (H2a–H4d), e também, separadamente, às comparações post hoc de Dunn. Para cada hipótese são relatados a estatística, o $p$ ajustado, o tamanho de efeito com IC e a margem de erro.
\textbf{Análise condicional (definida antes da coleta):} (i) se algum grupo de H2a ou H4b tiver $n<10$, a hipótese é relatada apenas de forma descritiva; (ii) se algum curso tiver $n<5$, os cursos são agrupados em bacharelados (CC, ES, SI) e tecnólogo (ADS), e H4a passa a usar Mann-Whitney; (iii) se o $\alpha$ de um construto for menor que 0,60, aplica-se a regra da Seção~5.3.
\textbf{Análises exploratórias (rotuladas como tal):} H4b repetida apenas com estudantes do 1º ao 4º período, para separar o efeito do estágio do avanço no curso; comparação por período (M8); comparação dos resultados entre canais de convite (sensibilidade à autosseleção).

\subsection{Respostas abertas}
ABR1 é analisada por codificação aberta indutiva (análise temática \citep{braun}). Dois integrantes codificam de forma independente 20\% das respostas, calcula-se o $\kappa$ de Cohen (com meta ${\geq 0{,}60}$), as divergências são resolvidas por consenso e o livro de códigos é consolidado. Em seguida, o restante é codificado e relata-se a frequência de cada tema (M15). Trechos que possam identificar alguém são removidos.

\subsection{Ferramentas}
A análise usa Python 3 (pandas, SciPy, pingouin), com scripts versionados no repositório. O dashboard da Parte 2 lê diretamente o CSV bruto exportado do formulário.

% =====================================================================
\Needspace{5\baselineskip}\section{Ética e LGPD}

\begin{itemize}[leftmargin=1em,itemsep=0.3ex]
\item \textbf{Enquadramento.} Trata-se de pesquisa de opinião com participantes não identificados, que a Resolução CNS nº 510/2016 (art. 1º, parágrafo único, I) dispensa de registro no sistema CEP/CONEP \citep{cns510}. Mesmo assim, o estudo segue as orientações da \pend{instituição} sobre pesquisas com pessoas e os princípios da resolução. Nenhuma lista de contatos de outras instituições é obtida: o convite circula apenas por intermediários e canais abertos.
\item \textbf{Consentimento.} A primeira página do formulário apresenta o TCLE (\texttt{instrumento/tcle.md}), que informa: objetivo, duração (cerca de 9 min), caráter voluntário, ausência de vínculo com notas, direito de desistir a qualquer momento antes do envio (após o envio, a resposta anônima não pode ser localizada), riscos mínimos (desconforto ao refletir sobre dificuldades) e contato dos responsáveis e do professor. A concordância é obrigatória para prosseguir (TRI1).
\item \textbf{Anonimização na origem.} O formulário não coleta nome, e-mail, matrícula, IP nem o nome da instituição; a idade é coletada em faixas. As respostas abertas são revisadas, e qualquer dado identificador é removido. O dashboard \emph{não exibe segmentos com menos de 5 respondentes}.
\item \textbf{LGPD.} Dados anonimizados não são dados pessoais (Lei nº 13.709/2018, art. 12) \citep{lgpd}. Ainda assim, adotam-se o consentimento (art. 7º, I) e os princípios de finalidade, necessidade e segurança (art. 6º). O uso é exclusivamente acadêmico.
\item \textbf{Armazenamento.} A exportação bruta fica na nuvem institucional, com acesso restrito ao grupo e ao professor, e será excluída logo após a defesa da Parte 2. O repositório público recebe apenas o conjunto anonimizado, sem textos livres não revisados.
\end{itemize}

% =====================================================================
\Needspace{5\baselineskip}\section{Ameaças à validade}

\begin{table}[H]
\centering\small
\caption{Ameaças típicas de survey e mitigação, na classificação de validade de Wohlin et al. (2012).}\label{tab:ameacas}
\begin{tabularx}{\textwidth}{@{}L{3.6cm}L{1.6cm}Y@{}}
\toprule
\thead{Ameaça} & \thead{Tipo} & \thead{Mitigação} \\
\midrule
Autosseleção e rede dos autores & Externa & Convite neutro (``sua experiência com programação''), sem a palavra ``dificuldade''; docentes e centros acadêmicos de outras instituições; registro do canal (PER5) e comparação entre canais. \\
Não resposta e cobertura & Externa & Dois lembretes; metas por segmento; comparação entre quem respondeu cedo e tarde \citep{armstrong}; distribuição por curso comparada ao Censo da Educação Superior (INEP). \\
Desejabilidade social & Construto & Anonimato; nenhum vínculo com notas; redação neutra; integrantes não avaliam os respondentes. \\
Aquiescência e descuido & Construto & Item invertido (DIF4); pergunta de atenção (ATN1); regras E5–E7. \\
Escalas adaptadas e traduzidas & Construto & Itens baseados em escalas validadas; pré-teste com especialista; $\alpha$ de Cronbach. A dificuldade \emph{percebida} não equivale a desempenho (limitação declarada). \\
Viés de memória & Construto & Instrução para considerar o aprendizado até agora; H4b repetida só com o 1º–4º período (exploratória). \\
Respostas automatizadas ou geradas por IA & Construto & Sem incentivos; tempo mínimo (E7); atenção (ATN1); duplicatas (E8); inspeção das respostas abertas. \\
Confusão e causalidade & Interna & Desenho transversal: apenas associações, em linguagem não causal; estágio e período analisados de forma estratificada (exploratória). \\
Generalização & Externa & Amostra não probabilística: conclusões restritas à amostra, margem nominal e composição descrita (cursos, instituições, canais). \\
Múltiplos testes, poder e $p$-hacking & Conclusão & Correção de Holm; poder declarado ($|\rho|\geq0{,}28$); hipóteses, testes e regras registrados antes da coleta. \\
Viés na codificação aberta & Conclusão & Dois codificadores; $\kappa$ de Cohen $\geq 0{,}60$; livro de códigos publicado. \\
\bottomrule
\end{tabularx}
\end{table}

% =====================================================================
\Needspace{5\baselineskip}\section{Plano de divulgação}

\textbf{Canais.} (C1) Docentes de programação de outras instituições, que recebem uma mensagem pronta para repassar às turmas e, quando autorizado, abrem 10 minutos da aula para a resposta, com QR code. (C2) Centros e diretórios acadêmicos e grupos de estudantes de tecnologia (WhatsApp, Discord, Telegram), com permissão dos administradores. (C3) Perfis dos integrantes no LinkedIn e no Instagram, com pedido de compartilhamento (bola de neve). (C4) Dois lembretes nos mesmos canais. Cada postagem é registrada com captura de tela datada, como evidência para a Parte 2.
\textbf{Monitoramento e contingência.} Em 25/11, a composição parcial é comparada às metas (15 respostas por curso e 20 estagiários), e a segunda onda é direcionada aos segmentos abaixo da meta. Se houver menos de 50 respostas válidas em 04/12, a coleta será estendida até 11/12, com novos canais, e a mudança será registrada como desvio.

\begin{table}[H]
\centering\small
\caption{Cronograma da coleta (2026).}\label{tab:crono}
\begin{tabularx}{\textwidth}{@{}L{2.9cm}Y@{}}
\toprule
\thead{Data} & \thead{Atividade} \\
\midrule
13/11 (sex) & Pré-teste, ajustes e registro do protocolo (commit e tag datados e envio na plataforma até 23:59) \\
16/11 e 17/11 & Contato com docentes e centros acadêmicos de outras instituições; preparação das peças de divulgação \\
18/11 (qua) & Abertura do formulário e primeira onda de divulgação (C1–C3) \\
25/11 (qua) & Verificação das metas por segmento; segunda onda direcionada e lembrete 1 (C4) \\
02/12 (qua) & Lembrete 2 (C4) \\
04/12 (sex) & Encerramento da coleta às 23:59; exportação bruta datada \\
a partir de 07/12 & Limpeza, análise, dashboard e registro de desvios (Parte 2) \\
\bottomrule
\end{tabularx}
\end{table}

% =====================================================================
{\renewcommand{\section}[2]{}%
\Needspace{6\baselineskip}\vspace{1.5ex}{\normalfont\bfseries\large Referências}\par\vspace{0.3ex}
\footnotesize\raggedright
\begin{multicols}{2}
\begin{thebibliography}{99}\setlength{\itemsep}{0pt}\setlength{\parskip}{0pt}
\bibitem[ARMSTRONG; OVERTON(1977)]{armstrong} ARMSTRONG, J. S.; OVERTON, T. S. Estimating nonresponse bias in mail surveys. \emph{Journal of Marketing Research}, v. 14, n. 3, p. 396--402, 1977.
\bibitem[BALTES; RALPH(2022)]{baltes} BALTES, S.; RALPH, P. Sampling in software engineering research: a critical review and guidelines. \emph{Empirical Software Engineering}, v. 27, 2022. DOI: 10.1007/s10664-021-10072-8.
\bibitem[BASILI; CALDIERA; ROMBACH(1994)]{basili} BASILI, V. R.; CALDIERA, G.; ROMBACH, H. D. The goal question metric approach. In: \emph{Encyclopedia of Software Engineering}. New York: Wiley, 1994.
\bibitem[BECKER et al.(2023)]{becker} BECKER, B. A. et al. Programming is hard -- or at least it used to be. In: \emph{Proc. 54th ACM Technical Symposium on Computer Science Education (SIGCSE)}. ACM, 2023. DOI: 10.1145/3545945.3569759.
\bibitem[BRAUN; CLARKE(2006)]{braun} BRAUN, V.; CLARKE, V. Using thematic analysis in psychology. \emph{Qualitative Research in Psychology}, v. 3, n. 2, p. 77--101, 2006.
\bibitem[BRASIL(2018)]{lgpd} BRASIL. Lei nº 13.709, de 14 de agosto de 2018. Lei Geral de Proteção de Dados Pessoais (LGPD). Brasília, 2018.
\bibitem[BRASIL(2016)]{cns510} BRASIL. Conselho Nacional de Saúde. Resolução nº 510, de 7 de abril de 2016. Brasília, 2016.
\bibitem[COHEN(1988)]{cohen} COHEN, J. \emph{Statistical power analysis for the behavioral sciences}. 2. ed. Hillsdale: Lawrence Erlbaum, 1988.
\bibitem[CONNOLLY; MURPHY; MOORE(2009)]{connolly} CONNOLLY, C.; MURPHY, E.; MOORE, S. Programming anxiety amongst computing students -- a key in the retention debate? \emph{IEEE Transactions on Education}, v. 52, n. 1, 2009. DOI: 10.1109/TE.2008.917193.
\bibitem[HOLM(1979)]{holm} HOLM, S. A simple sequentially rejective multiple test procedure. \emph{Scandinavian Journal of Statistics}, v. 6, n. 2, p. 65--70, 1979.
\bibitem[KITCHENHAM; PFLEEGER(2008)]{kitchenham} KITCHENHAM, B. A.; PFLEEGER, S. L. Personal opinion surveys. In: SHULL, F.; SINGER, J.; SJØBERG, D. (ed.). \emph{Guide to Advanced Empirical Software Engineering}. London: Springer, 2008. DOI: 10.1007/978-1-84800-044-5\_3.
\bibitem[LAHTINEN; ALA-MUTKA; JÄRVINEN(2005)]{lahtinen} LAHTINEN, E.; ALA-MUTKA, K.; JÄRVINEN, H.-M. A study of the difficulties of novice programmers. In: \emph{Proc. 10th ITiCSE}. ACM, 2005. DOI: 10.1145/1067445.1067453.
\bibitem[RALPH et al.(2021)]{ralph} RALPH, P. et al. \emph{ACM SIGSOFT Empirical Standards}: Questionnaire Surveys; Registered Reports. 2021. Disponível em: https://github.com/acmsigsoft/EmpiricalStandards.
\bibitem[RAMALINGAM; WIEDENBECK(1998)]{ramalingam} RAMALINGAM, V.; WIEDENBECK, S. Development and validation of scores on a computer programming self-efficacy scale and group analyses of novice programmer self-efficacy. \emph{Journal of Educational Computing Research}, v. 19, n. 4, 1998. DOI: 10.2190/C670-Y3C8-LTJ1-CT3P.
\bibitem[ROBINS; ROUNTREE; ROUNTREE(2003)]{robins} ROBINS, A.; ROUNTREE, J.; ROUNTREE, N. Learning and teaching programming: a review and discussion. \emph{Computer Science Education}, v. 13, n. 2, p. 137--172, 2003.
\bibitem[ROMANO et al.(2006)]{romano} ROMANO, J. et al. Appropriate statistics for ordinal level data. In: \emph{Annual Meeting of the Florida Association of Institutional Research}, 2006.
\bibitem[TOMCZAK; TOMCZAK(2014)]{tomczak} TOMCZAK, M.; TOMCZAK, E. The need to report effect size estimates revisited. \emph{Trends in Sport Sciences}, v. 1, n. 21, p. 19--25, 2014.
\bibitem[TSAI; WANG; HSU(2018)]{tsai} TSAI, M.-J.; WANG, C.-Y.; HSU, P.-F. Developing the computer programming self-efficacy scale for computer literacy education. \emph{Journal of Educational Computing Research}, 2018. DOI: 10.1177/0735633117746747.
\bibitem[WATSON; LI(2014)]{watson} WATSON, C.; LI, F. W. B. Failure rates in introductory programming revisited. In: \emph{Proc. ITiCSE '14}. ACM, 2014. DOI: 10.1145/2591708.2591749.
\bibitem[WOHLIN et al.(2012)]{wohlin} WOHLIN, C. et al. \emph{Experimentation in Software Engineering}. Berlin: Springer, 2012.
\end{thebibliography}
\end{multicols}}

\end{document}
